\section{Residential air-pollutant results}
\label{sec:air}
The survey-based S reference produces \SI{\SectorRefPMfine}{\kilo\tonne} primary PM$_{2.5}$ and \SI{\SectorRefBC}{\kilo\tonne} BC in FY2023--24. The latest common-core year is FY2024--25. Air mass and GHG results use the same household fuel configurations, while species-specific factors retain their own units and process definitions.

\begin{table}[htbp]\centering\small
% Native Gg; missing species are NE, not zero.
\begin{tabular}{@{}lrr@{}}\toprule
Species (Gg) & FY2023--24 & FY2024--25\\\midrule
PM$_{2.5}$ & 282.8 & 276.4 \\
BC & 20 & 19.55 \\
NO$_x$ & 156.1 & 154.2 \\
SO$_2$ & 20.84 & 20.8 \\
CO & 3354 & 3265 \\
NMVOC & 357.7 & 357.9 \\
NH$_3$ & 42.38 & 41.21 \\
\bottomrule\end{tabular}

\caption{Gross appliance emissions under reference S. BC and OC are components of PM$_{2.5}$; PM$_{10}$/TSP have narrower coal-only support.}\label{tab:air}
\end{table}

\subsection{Factors and material boundary}
Andreae biomass factors use dry fuel including ash. Dung uses total dry solids rather than ash-free volatile solids, wood is converted from survey air-dry quantity, and household residue burning is separate from field burning. Fossil gas, LPG, kerosene and coal use applicable EMEP small-combustion factors on energy. LPG's gaseous-appliance proxy and ambiguous filterable/condensable definitions remain explicit. Household generators and unusable charcoal sulfur factors remain missing estimates.

NO$_x$ is NO$_2$-equivalent throughout. Original NO factors receive the 46/30 conversion once. Pollutants are native masses and do not enter a CO$_2$-equivalent aggregate. Nested particulate species are not summed.

\subsection{Configurations and uncertainty}
S uses household dung reported by the HIES reconstruction. L assumes that all agricultural-ledger diverted dung is burned in households; the transfer establishes available material, not its destination. The alternative changes dung emissions consistently in GHG and air products, without assigning a probability to that fate. In the national FY2023--24 air panel, this case changes covered PM$_{2.5}$ from about 423 to 772 Gg.

The conditional ensemble shares household activity across pollutants. Ledger-L dung uses the same sampled herd as manure accounting, with diversion and volatile-solid assumptions held. Stove/fuel factor transfer and structural fate cases are separate from parameter intervals. The former road-factor technology endpoints are sensitivities, not a probabilistic basis for a residential/road ratio interval.

\subsection{Fitness for use}
Residential biomass is the largest covered PM$_{2.5}$ source under S, but pollutant and GHG rankings differ by specific process; there is no general inverse ordering. The source comparison is gross emitted mass, not exposure or a health-intervention ranking.

The product reports emissions at the appliance. Indoor deposition and outdoor escape are not estimated, so this mass is not an ambient-ready outdoor release. Stove location, ventilation and personal time--activity patterns would require separate evidence. Survey national/provincial controls and fixed geographic proxies support allocations rather than observed changes in local combustion. Monthly source evidence and held spatial patterns are labelled accordingly.
